% Nová příprava: 2026-10-28; 1. ročník SPŠ, etapa 13–24.
% Rozpis: 1_rocnik/tematicke_plany/Rozpis_vyuky_MAT_1_rocnik.tex, hodina 20.
% Původ: vlastní formulace výkladu a vlastní úlohy; metodická inspirace uvedenými zdroji.
% Kontrola, obsahové porovnání zdrojů a časy: planovani/DOKONCENI_PRIPRAV_13_24.md.
% Učebnice: ucebnice/01_Zakladni_poznatky/Mnoziny/Množiny.tex
% oddíly Rovnost množin, Podmnožiny, Množinové operace a Vennovy diagramy; dle tématu hodiny.
% Realisticky.cz: externi_zdroje/realisticky/01_Zakladni_poznatky/03_Mnoziny/010304_Vennovy_diagramy/010304_Lekce.pdf
% Realisticky.cz: externi_zdroje/realisticky/01_Zakladni_poznatky/03_Mnoziny/010308_Intervaly/010308_Lekce.pdf
% Realisticky.cz: externi_zdroje/realisticky/01_Zakladni_poznatky/03_Mnoziny/010309_Mnoziny_shrnuti/010309_Lekce.pdf
\documentclass[11pt,a4paper]{article}
% Spolecny styl pro ucitelske pripravy MAT (Work/MAT)
% Verze 1.1, 2026-10-23
\NeedsTeXFormat{LaTeX2e}
\ProvidesPackage{priprava}[2026/10/23 v1.1 Sablona ucitelskych priprav MAT]

% Zakladni typografie a cestina
\RequirePackage[T1]{fontenc}
\RequirePackage[utf8]{inputenc}
\RequirePackage[czech]{babel}
\RequirePackage{lmodern}
\RequirePackage{microtype}
\RequirePackage{amsmath,amssymb}
\RequirePackage{xcolor}
\RequirePackage{enumitem}
\RequirePackage[most]{tcolorbox}
\RequirePackage{tabularx,booktabs,array}
\RequirePackage{fancyhdr}
\RequirePackage{needspace}
\RequirePackage{geometry}
\RequirePackage[hidelinks]{hyperref}

\geometry{a4paper,top=17mm,bottom=18mm,left=23mm,right=23mm}
\definecolor{notegray}{gray}{0.39}
\setlength{\parindent}{0pt}
\setlength{\parskip}{2.5pt plus 1pt minus 0.5pt}
\setlist[itemize]{topsep=2pt,itemsep=1pt,parsep=0pt,leftmargin=1.45em}
\setlist[enumerate]{topsep=2pt,itemsep=2pt,parsep=0pt,leftmargin=1.65em}
\setcounter{secnumdepth}{0}
\renewcommand{\arraystretch}{1.12}

% Zahlavi: vlevo rocnik | globalni tema, vpravo cislo: nazev hodiny
% Zapati zustava zcela prazdne.
\pagestyle{fancy}
\fancyhf{}
\renewcommand{\headrulewidth}{0.35pt}
\renewcommand{\footrulewidth}{0pt}
\setlength{\headheight}{14pt}
\newcommand{\nastavPripravu}[4]{%
  \fancyhead[L]{\small #1\enspace\textbar\enspace #2}%
  \fancyhead[R]{\small #3: #4}%
}

% Minimalisticke cernobile boxy. Kazda definice zvlast.
\tcbset{commonlesson/.style={enhanced,breakable=false,colback=white,colframe=black,
  colbacktitle=white,coltitle=black,boxrule=0.45pt,
  arc=0.8mm,outer arc=0.8mm,left=3mm,right=3mm,top=1.6mm,bottom=1.6mm,
  toptitle=1mm,bottomtitle=0.7mm,fonttitle=\bfseries,
  attach boxed title to top left={xshift=2.5mm,yshift=-1.15mm},
  boxed title style={colback=white,colframe=white,boxrule=0pt,left=1mm,right=1mm}}}
\newtcolorbox{zadaniBox}[1][]{commonlesson,title={Zadání},#1}
\newtcolorbox{prikladBox}[1][]{commonlesson,title={Příklad},#1}
\newtcolorbox{definiceBox}[1][]{commonlesson,title={Definice},#1}

% Spolecny box pro zadani a jeho reseni, bez mezery mezi castmi.
% Pouziti:
%   \begin{prikladReseniBox}
%     Zadani ulohy.
%     \tcblower
%     \begin{align*} ... \end{align*}
%   \end{prikladReseniBox}
% Spodni cast ma jemny sedy podklad a delici linku; lze v ni sazet
% vice radku matematiky. Obe casti maji jen jeden spolecny ramecek.
\tcbset{solutionlesson/.style={
  commonlesson,
  enhanced jigsaw,
  breakable,
  colbacklower=black!3,
  segmentation style={solid,draw=black!25,line width=0.35pt},
  fontlower=\small\color{notegray},
  before lower={\noindent\textbf{Řešení:}\enspace},
  leftlower=3mm,rightlower=3mm
}}
\newtcolorbox{prikladReseniBox}[1][]{solutionlesson,title={Příklad},#1}
\newtcolorbox{zadaniReseniBox}[1][]{solutionlesson,title={Zadání},#1}

% Pomucky pro autora pripravy: sedive poznamky; \reseni pro starsi pripravy
\newcommand{\poznamka}[1]{\par\begingroup\small\color{notegray}\itshape
  \textnormal{\textbf{Poznámka pro učitele.}}\enspace #1\par\endgroup}
\newcommand{\reseni}[1]{\par\begingroup\small\color{notegray}
  \textbf{Očekávané řešení:}\enspace #1\par\endgroup}
\newcommand{\blok}[1]{\needspace{5\baselineskip}\subsection*{#1}}
\newcommand{\vyrok}[1]{\uv{#1}}

\nastavPripravu{1. ročník}{Množiny}{20}{Opakování množin}
\begin{document}
\blok{Nejdříve zvol význam, potom počítej}
Nech žáky bez sešitu vysvětlit rozdíly mezi $\in$ a $\subseteq$,
$\varnothing$ a $\{0\}$, $\cap$ a $\cup$, rozdílem a doplňkem.
Podle odpovědí zvol krátké připomenutí; dnes nezavádíme nový povinný pojem.
Shrnutí pravidel má pomoci řešit nový příklad, nikoli nahradit jeho rozbor.

\begin{prikladReseniBox}[title={P1: tři modelové chyby},breakable=false]
Oprav: a) $\{0\}=\varnothing$; b) $\{1,2\}\subseteq\{2,3\}$;
c) $(-1;2)\cup(2;4)=(-1;4)$.
\tcblower
a) Nula je prvkem $\{0\}$, prázdná množina žádný prvek nemá.
b) Vztah neplatí kvůli prvku $1$.
c) Bod $2$ nepatří do žádného původního intervalu. Výsledek musí zůstat
$(-1;2)\cup(2;4)$; krajní hodnoty samy nevystihnou chybějící bod.
\end{prikladReseniBox}

\begin{prikladReseniBox}[title={P2: složitější výraz rozlož na kroky},breakable=false]
V $U=\{-3,-2,-1,0,1,2,3,4\}$ je $A$ množina sudých čísel,
$B=\{x\in U\mid |x|\le2\}$ a $C=\{x\in U\mid x>0\}$.
Urči $(A\cup B)\setminus C$, $(A\setminus B)\cap C$ a $(B\cap C)'_U$.
\tcblower
$A=\{-2,0,2,4\}$, $B=\{-2,-1,0,1,2\}$, $C=\{1,2,3,4\}$.
\begin{align*}
(A\cup B)\setminus C
 &=\{-2,-1,0,1,2,4\}\setminus\{1,2,3,4\}=\{-2,-1,0\},\\
(A\setminus B)\cap C&=\{4\}\cap\{1,2,3,4\}=\{4\},\\
(B\cap C)'_U&=\{1,2\}'_U=\{-3,-2,-1,0,3,4\}.
\end{align*}
Každý výsledek je podmnožinou $U$. Doplněk posledního mezivýsledku
neznamená „opačná čísla“.
\end{prikladReseniBox}
\poznamka{Jádro: úvodní vybavení pojmů, P1, P2 a S1 a)–d),
asi 38 min práce, organizace a závěr 7 min. S1 e) je navazující úloha
pro rychlé, R1 další rezerva. U každého typu požaduj zdůvodnění,
ne pouze značku nebo výsledek.}

\blok{Samostatné procvičení napříč tématy}
\begin{zadaniReseniBox}[title={S1: samostatný pokus a společná kontrola},breakable=false]
\begin{enumerate}[label=\alph*)]
\item Vypiš $M=\{x\in\mathbb Z\mid-2\le x<2\}$.
Platí $M\subseteq\langle-2;2)$ a platí obrácený vztah?
\item V $U=\{1,2,3,4,5,6\}$ pro $A=\{1,2,3\}$, $B=\{3,4,5\}$
urči $A\cap B$, $A\cup B$, $A\setminus B$, $(A\cup B)'_U$.
Umísti prvky do čtyř oblastí diagramu.
\item Pro $I=\langle-4;1)$ a $J=\langle-1;3\rangle$ urči $I\cap J$, $I\cup J$.
\item Z 27 výrobků má vlastnost $A$ 14, vlastnost $B$ 12 a žádnou z nich 5 výrobků.
Urči počet s oběma vlastnostmi a s právě jednou.
\item V $U=\{1,2,3,4\}$, $A=\{1,2\}$, $B=\{2,3\}$ a $C=\{3,4\}$
urči $A'\cap(B\cup C)$.
\end{enumerate}
\tcblower
a) $M=\{-2,-1,0,1\}$, první zahrnutí ano, druhé ne (např. $\frac12$).
b) $\{3\}$, $\{1,2,3,4,5\}$, $\{1,2\}$, $\{6\}$.
Jen $A$: $1,2$; obě: $3$; jen $B$: $4,5$; mimo obě: $6$.
c) $\langle-1;1)$ a $\langle-4;3\rangle$.
d) Sjednocení $27-5=22$, obě $14+12-22=4$.
Právě jedna $(14-4)+(12-4)=18$; kontrola $10+4+8+5=27$.
e) $A'=\{3,4\}$, $B\cup C=\{2,3,4\}$, průnik $\{3,4\}$.
\end{zadaniReseniBox}

\blok{Volitelná rezerva}
\begin{zadaniReseniBox}[title={R1: od příkladu k obecnému zdůvodnění},breakable=false]
\begin{enumerate}[label=\alph*)]
\item Dokaž, že $A\cup B=A\cap B$ platí právě tehdy, když $A=B$.
\item Urči doplněk $\langle-2;1)\cup(1;4\rangle$ v $\mathbb R$.
\end{enumerate}
\tcblower
a) Je-li $A=B$, obě operace vrátí tutéž množinu.
Platí-li rovnost operací, každý prvek $A$ je ve sjednocení, tedy také
v průniku a tím v $B$; stejnou úvahou $B\subseteq A$.
Obě zahrnutí dávají $A=B$.
b) $(-\infty;-2)\cup\{1\}\cup(4;+\infty)$.
Nezapomeň na vnitřní chybějící bod $1$; $-2$ a $4$ do původní množiny patří.
\end{zadaniReseniBox}
\poznamka{Rezerva 10–15 min. Závěrem ať každý žák uvede vlastní kontrolní otázku,
např. „Ověřil jsem oba krajní body?“ a předvede ji na konkrétní úloze.
Zjištěné potíže využij pro rozbor v hodině 21.}

\blok{Domácí procvičení před průběžným testem}
\begin{zadaniReseniBox}[title={D1: 12–18 minut}]
\begin{enumerate}[label=\alph*)]
\item Vypiš všechny podmnožiny $\{0,5\}$.
\item V $U=\{0,1,2,3,4,5\}$ je $K=\{0,2,4\}$ a $L=\{2,3,5\}$.
Urči $K\cap L$, $K\cup L$, $L\setminus K$ a $K'_U$.
\item Pro $I=(-3;2\rangle$ a $J=(2;6)$ urči průnik, sjednocení a ověř bod $2$.
\item Z 24 lidí používá nástroj $A$ 13, nástroj $B$ 10 a oba 4 lidé.
Kolik nepoužívá ani jeden? Nakresli diagram s počty.
\item Platí vždy $(A\cup B)'=A'\cup B'$? Uveď protipříklad se stejným $U$.
\end{enumerate}
\tcblower
a) $\varnothing,\{0\},\{5\},\{0,5\}$.
b) $\{2\}$, $\{0,2,3,4,5\}$, $\{3,5\}$, $\{1,3,5\}$.
c) $\varnothing$, $(-3;6)$. Dvojka je ve sjednocení díky $I$, do průniku nepatří.
d) Jen $A$: $9$, obě $4$, jen $B$: $6$, žádný $24-(13+10-4)=5$.
e) Ne. Pro $U=\{1,2\}$, $A=\{1\}$, $B=\{2\}$ je vlevo
$\varnothing$ a vpravo $U$. Správně je průnik doplňků.
\end{zadaniReseniBox}
\end{document}
